% Lösungen Einheit 1 von 5 – Rechteck, Quadrat, Umfang (zusammenbau v0.1)
\einheitenkopf{Einheit 1 von 5 · Rechteck, Quadrat, Umfang}

\erg{10}{a) $a$ einkreisen – die Länge ist gesucht}
\erg{11}{a) Umfang \quad b) $A = 8 \cdot 3 = 24$ cm², $u = 2 \cdot (8 + 3) = 22$ cm \quad c) $A = 4{,}5 \cdot 2 = 9$ cm², $u = 2 \cdot (4{,}5 + 2) = 13$ cm \quad d) $A = 6 \cdot 6 = 36$ cm², $u = 4 \cdot 6 = 24$ cm \quad e) $u = 3 + 4 + 5 + 2 + 6 = 20$ cm \quad f) $b = 56 : 8 = 7$ cm \quad g) $b = (20 - 2 \cdot 6) : 2 = 4$ cm \quad h) $a = \sqrt{64} = 8$ cm \quad i) $u = 6x$ (zweimal $2x$ und zweimal $x$) \quad j) $34 : 2 = 17$; $17 - 11 = 6$ cm}
\erg{12}{a) $24 + 6 = 30$ m²}
\erg{13}{a) $2 \cdot 0{,}45 = 0{,}9$ m²}
\erg{14}{a) Tom hat die Fläche berechnet, nicht den Umfang. Richtig: $u = 2 \cdot (6 + 2{,}5) = 17$ cm \quad b) In eine Reihe passen so viele Einheitsquadrate, wie das Rechteck lang ist; es gibt so viele Reihen, wie es breit ist – also Länge mal Breite. \quad c) $A = 4{,}2 \cdot 3{,}5 = 14{,}7$ m²; $14{,}7 : 2{,}2 \approx 6{,}7$, also $7$ Pakete \quad d) Rechteck mit den Seiten $3$ und $x$}

14d)

\viereck[seiten={$x$,$3$,$x$,$3$}]{(0,0)}{(4,0)}{(4,2)}{(0,2)}

